\section*{Zahvalnica}
Zahvaljujem se...

\newpage

\section*{Rezime}

Ubrzani razvoj veštačke inteligencije i velikih jezičkih modela (engl.
\emph{large language models}, LLM) menja praksu provere znanja u visokom obrazovanju,
naročito u ocenjivanju pitanja otvorenog tipa.
Tradicionalno ocenjivanje je vremenski zahtevno, sklono nedoslednosti i teško skalabilno za velike grupe studenata.
Veliki jezički modeli tumače slobodan tekst i rezonuju o njegovom značenju,
pa povratna informacija može biti brža, doslednija i češća.
U ovom radu predstavljen je EduGrader, platforma za ocenjivanje otvorenog koda koja objedinjuje GPT-4o,
GPT-4o-mini, DeepSeek-Chat, GPT-5 i DeepSeek-Reasoner.
EduGrader podržava tri konfiguracije strogosti ocenjivanja (blagu, neutralnu,
strogu) i uz numeričku ocenu daje sažeto obrazloženje: šta je u odgovoru ispravno, šta nedostaje i gde je greška.
Sistem radi u dva režima: režimu sa zadatim referentnim odgovorom (engl. \emph{provided-reference-answer}, PRA),
vođenom referentnim odgovorima nastavnika,
i režimu sa generisanim referentnim odgovorom (engl. \emph{generated-reference-answer}, GRA),
koji referentne odgovore automatski konstruiše iz nastavnih materijala.
Svi eksperimenti sprovedeni su na srpskom, jeziku sa ograničenim resursima koji je slabo zastupljen u podacima za obučavanje savremenih modela,
dakle izvan engleskog konteksta u kome se ovakvi sistemi obično ispituju.
EduGrader je evaluiran na 583 autentična studentska odgovora sa pet univerzitetskih kurseva na dve institucije,
uključujući tekstualna objašnjenja i odgovore u obliku koda.
U najboljoj konfiguraciji sistem postiže Pirsonov koeficijent korelacije od 0{,}89 sa ljudskim ocenjivačima.
Zamišljen je kao pomoć nastavniku koja preuzima rutinski deo ocenjivanja, a ne kao njegova zamena.

\vspace{1em}
\noindent\textbf{Ključne reči:} veliki jezički modeli, automatsko ocenjivanje, pitanja otvorenog tipa, referentni odgovor,
srpski jezik, jezik sa ograničenim resursima, procena znanja u visokom obrazovanju, EduGrader.
\newpage
